\documentclass[10pt,a4paper]{amsart}
\usepackage[margin=19mm]{geometry}
\usepackage{amsmath,amssymb,mathtools}
\usepackage[scr=rsfs,cal=euler]{mathalfa}
\usepackage{xfrac}
\usepackage[unicode,hidelinks,bookmarks=false]{hyperref}
\newcommand{\ie}{i.e.\ }
\newcommand{\cf}{cf.\ }
\newcommand{\resp}{respectively\ }
\newcommand{\Subsection}{\subsection}
\providecommand{\nobreakdash}{\nobreak\hbox{-}}
\setlength{\emergencystretch}{2em}
\multlinegap0pt
\usepackage{amssymb}
\newtheorem{lem}{Lemma}
\newcommand{\Comp}{\mathbb C}
\newcommand{\eps}{\varepsilon}
\newcommand{\norm}[1]{\| #1\|}
\newcommand{\tr}{\mathrm{T}}
\title{On the small boundary property and $\mathcal Z$-absorption}
\author{George A. Elliott, Zhuang Niu}
\date{Source: arXiv:2504.03611v3 (CC BY 4.0)}
\begin{document}
\maketitle

\noindent\emph{Source and licence.} On the small boundary property and $\mathcal Z$-absorption. Version arXiv:2504.03611v3 (CC BY 4.0), distributed by arXiv under the Creative Commons Attribution 4.0 International licence, http://creativecommons.org/licenses/by/4.0/, as recorded in the arXiv licence metadata for this exact version and shown on the abstract page https://arxiv.org/abs/2504.03611v3. Full licence notice: \texttt{licenses/arxiv-2504.03611-CC-BY-4.0.txt}.\par\smallskip

\noindent\textit{Editorial scope.} Counted unit: the article's lem D-sandwich (source0.tex lines 1492--1506). The article's complete original proof of it is delivered with it (source0.tex lines 1508--1660, one \texttt{proof} environment of 153 lines). No mathematics is written, repaired or re-derived: the statement and the proof are the article's own lines, in the article's own notation, and no retained line is substituted or re-flowed.  No further source-local context is needed: every cross-reference of every retained line resolves inside the two delivered spans.  Nothing was omitted from any retained span, so the package carries no omission record.  No retained line contains a citation, so the wrapper carries no bibliography and no bibliography entry.  No retained line contains a figure environment, an image-inclusion command or drawing code, so no image is delivered and the package declares no asset dependency.  Editorial formatting only, in addition: an AMS \texttt{amsart} preamble that restates the article's own declarations and macros used by the retained lines, namely the environments \texttt{lem} on separate counters, together with the standard packages those lines need.  The retained environments print in this document as Lemma 1 in order of appearance, whereas the article prints the same items with the numbering of its own full document.  The complete native source of the article as distributed in the arXiv e-print for this exact version is bundled unchanged as \texttt{sources/175995/source0.tex}.
\begin{lem}\label{D-sandwich}
Let $A$ be a unital C*-algebra, and let $D \subseteq A$ be a unital commutative subalgebra such that the pair $(D, A)$ has Property (C). Let $\phi_1, \phi_2, \phi_3, \psi_1, \psi_2 \in (D)_1^+$ and $\eps_0>0$ have the following properties:
\begin{enumerate}
\item\label{cond-s-1} $\phi_2, \phi_3 \in \overline{\phi_1 D\phi_1}$,
\item\label{cond-s-2} $\psi_1\psi_2 = \psi_2$, $\phi_2 \phi_3 = \phi_3$,
\item\label{cond-s-3} $\mathrm{d}_\tau(\psi_1) < \mathrm{d}_\tau(\phi_1)$ for all $\tau\in\tr(A)$,
\item\label{cond-s-4} $\inf\{\tau(\psi_2) - \mathrm{d}_\tau(\phi_3): \tau \in \tr(A)\} >0 $, and
\item\label{cond-s-5} $\inf\{\mathrm{d}_\tau(\phi_2) - \tau(\eta_{\eps_0}(\psi_1)): \tau \in \tr(A)\} >0 $.  \end{enumerate}
Then, for any $\eps>0$, there is a contraction $u\in A$ such that
$$ u^*\psi_1u \in_\eps^{\norm{\cdot}_2}  \overline{\phi_1A\phi_1},\quad u^*\eta_{\eps_0/2}(\psi_1)u \in^{\norm{\cdot}_2}_{\eps} \overline{\phi_2A\phi_2}, \quad  \eta_{\eps}(u^*\psi_1u) \phi_3 \approx_{\eps}^{\norm{\cdot}_2} \phi_3,$$
and
$$\mathrm{dist}_{2, \tr(A)}(udu^*, (D)_1) < \eps,\quad \mathrm{dist}_{2, \tr(A)}(u^*du, (D)_1) < \eps,\quad d \in (D)_1,$$
and
$$\norm{uu^* - 1}_{2, \tr(A)}, \quad \norm{u^*u - 1}_{2, \tr(A)} < \eps.$$
\end{lem}

\begin{proof}
Let $\eps>0$ be given. Choose $\delta \in (0, \min\{\eps_0, \eps\}/4)$ such that if $x, y$ are positive contractions of a unital C*-algebra $A$ such that $\norm{x - y}_{2, \tr(A)} < \delta$, then
\begin{equation}\label{pre-pert-delta}
\norm{\eta_\eps(x) - \eta_\eps(y)}_{2, \tr(A)} < \eps/2.
\end{equation}


Define
$$ \delta_1 := \inf\{\tau(\psi_2) - \mathrm{d}_\tau(\phi_3): \tau \in \tr(A)\} >0$$
and
$$\delta_2 := \inf\{\mathrm{d}_\tau(\phi_2) - \tau(\eta_{\eps_0}(\psi_1)): \tau \in \tr(A)\} > 0. $$

By Property (C) and Assumption (\ref{cond-s-3}), for any $\eps'>0$ (to be determined later), there is a contraction $u_1\in A$ such that
\begin{equation}\label{approx-C-1-h}
u^*_1 \psi_1 u_1 \in^{\norm{\cdot}_2}_{\eps'} \overline{\phi_1 A \phi_1}, 
\end{equation}
\begin{equation}\label{approx-C-1-n}
\mathrm{dist}_{2, \tr(A)}(u_1du^*_1, (D)_1) < \eps', \quad \mathrm{dist}_{2, \tr(A)}(u_1^*du_1, (D)_1) < \eps',\quad d \in (D)_1, 
\end{equation}
and
\begin{equation}\label{approx-C-1-u}
\norm{u_1u_1^* - 1}_{2, \tr(A)},\quad  \norm{u_1^*u_1 - 1}_{2, \tr(A)} < \eps'.
\end{equation}

By \eqref{approx-C-1-h}, there is $n$ large enough that $$ \norm{(\phi_1^{\frac{1}{n}}) (u_1^*\psi_1u_1) - u_1^*\psi_1u_1}_{2, \tr(A)} < \eps'. 
$$
By \eqref{approx-C-1-n}, there is a positive contraction $d_1 \in D$ such that
\begin{equation*}
\norm{d_1- u_1^*\psi_1u_1}_{2, \tr(A)}< \eps', 
\end{equation*}
and then
$$\norm{\phi_1^{\frac{1}{n}}d_1 - u_1^*\psi_1u_1}_{2, \tr(A)} < 2\eps'. 
$$
Consider $\phi_1^{\frac{1}{n}}d_1$, and still denote this element by $d_1$. One has $$d_1 \in \overline{\phi_1 D \phi_1}$$
and
\begin{equation}\label{approx-C-1}
\norm{d_1 - u_1^*\psi_1u_1}_{2, \tr(A)} < 2\eps'. 
\end{equation}


Consider  the contraction
$$\eta_{\eps_0}(u^*_1\psi_1 u_1),$$
and note that (by \eqref{approx-C-1}, \eqref{approx-C-1-u} and Condition \eqref{cond-s-5})
$$ \mathrm{d}_\tau( \eta_{\eps_0/2}(d_1) ) \leq \tau(\eta_{\eps_0}(d_1))  \approx_{O(\eps')}^{\norm{\cdot}_2}  \tau( \eta_{\eps_0}(u^*_1\psi_1 u_1) ) \approx_{O(\eps')}^{\norm{\cdot}_2} \tau( \eta_{\eps_0}(\psi_1) ) < \mathrm{d}_\tau(\phi_2) - \delta_2/2,\quad \tau \in \tr(A).$$
With $\eps'$ sufficiently small, one has 
$$\mathrm{d}_\tau(\eta_{\eps_0/2}(d_1)) < \mathrm{d}_\tau(\phi_2),\quad \tau \in \tr(A).$$ 
By Property (C), for any $\eps''>0$ (to be determined later), there is $u_2 \in \overline{\phi_1 A \phi_1}+\Comp 1$ such that 
\begin{equation}\label{approx-C-2-here}
u^*_2 \eta_{\eps_0/2}(d_1) u_2 \in^{\norm{\cdot}_2}_{\eps''} \overline{\phi_2 A \phi_2},
\end{equation}
\begin{equation}\label{approx-C-2-n}
\mathrm{dist}_{2, \tr(A)}(u_2 d u^*_2, D_1) < \eps'', \quad \mathrm{dist}_{2, \tr(A)}(u_2^*du_2, D_1) < \eps'',\quad d \in D_1,
\end{equation}
and
\begin{equation}\label{approx-C-2-u}
\norm{u_2u_2^* - 1}_{2, \tr(A)},\quad  \norm{u_2^*u_2 - 1}_{2, \tr(A)} < \eps''.
\end{equation}


Since (as $\delta < \eps_0/4$)
$$\eta_{\eps_0/2}(d_1) \eta_\delta(d_1) = \eta_\delta(d_1),$$
it follows from \eqref{approx-C-2-here} that
\begin{equation}\label{close-hereditary}
u_2^* \eta_\delta(d_1) u_2 \in^{\norm{\cdot}_2}_{2\eps''} \overline{\phi_2 A \phi_2}. 
\end{equation}

By \eqref{approx-C-2-n}, there are positive contractions $$ d_{1, 1},\  d_{1, \delta} \in  D $$ such that
\begin{equation}\label{pert-C-step-2}
 \norm{d_{1, 1} - u_2^* d_1 u_2}_{2, \tr(A)} < \eps'',\quad  \norm{d_{1, \delta} - u_2^*\eta_{\delta}(d_1)u_2}_{2, \tr(A)} < \eps''. 
\end{equation}  
 By \eqref{close-hereditary}, $$d_{1, \delta} \in_{3\eps''}^{\norm{\cdot}_2} \in \overline{\phi_2A\phi_2}.$$
 With the same argument as above for $d_1$, there is a positive contraction, still denoted by $ d_{1, \delta}$, such that 
 \begin{equation}\label{u-2-2-approx-C-1}
 d_{1, \delta} \in \overline{\phi_2 D \phi_2}\quad \mathrm{and}\quad \norm{d_{1, \delta} - u_2^*\eta_{\delta}(d_1)u_2}_{2, \tr(A)} < 5\eps''.
 \end{equation}
Also note that, by \eqref{approx-C-2-u}, $$ \norm{ \eta_\delta(d_{1, 1}) - d_{1, \delta} }_{2, \tr(A)} = O(\eps''). $$

Define
$$\tilde{d}_{1, 1} = f_{\delta}(d_{1, 1}) \quad \mathrm{and} \quad \tilde{d}_{1, \delta} = \eta_{\delta}(d_{1, 1}),$$
where
$$
f_{\delta}(t) = \left\{ 
\begin{array}{ll}
0, & t \leq 0, \\
\mathrm{linear}, & t\in [0, 1-\delta], \\
1, & t \geq 1-\delta.
\end{array}
\right.
$$ 
Then
\begin{equation}\label{pert-fur-hereditary}
 \tilde{d}_{1, 1} \tilde{d}_{1, \delta} = \tilde{d}_{1, \delta} \quad \mathrm{and} \quad \norm{\tilde{d}_{1, 1} - d_{1, 1}} < \delta,
 \end{equation}
and by \eqref{pert-C-step-2},
$$ \norm{\tilde{d}_{1, \delta} - d_{1, \delta}}_{2, \tr(A)} \approx_{O(\eps'')} \norm{\eta_\delta(u_2^* d_1 u_2) -  u^*_2\eta_\delta(d_1)u_2} = O(\eps'').$$
Then, with $\eps''$ sufficiently small, there is $n \in \mathbb N$ such that $$ \norm{(\tilde{d}_{1, \delta})^{\frac{1}{n}} d_{1, \delta} - d_{1, \delta}}_{2, \tr(A)} < \delta_1/4, $$
and hence, for all $\tau \in \tr(A)$,
\begin{eqnarray*}
\mathrm{d}_\tau((\tilde{d}_{1, \delta})^{\frac{1}{n}} d_{1, \delta}) + \delta_1/4 & \geq & \tau((\tilde{d}_{1, \delta})^{\frac{1}{n}} d_{1, \delta}) + \delta_1/4   \\
%& > & \tau(d_{1, \delta}) \\
%& > & \mathrm{d}_\tau(\tilde{d}_{1, \delta})   >  \tau(\tilde{d}_{1, \delta}) \\
& > & \tau(d_{1, \delta}) \approx_{5\eps''} \tau(u_2^* \eta_{\delta}(d_1)u_2) \quad \quad (\eqref{u-2-2-approx-C-1}) \\
& \approx_{O(\eps')} &  \tau(u^*_2\eta_{\delta}(u_1^* \psi_1 u_1 )u_2) \quad\quad (\eqref{approx-C-1}) \\
& \approx_{O(\eps' + \eps'')} & \tau(\eta_{\delta}(\psi_1 ))  > \tau(\psi_2) \\ 
& > & \mathrm{d}_\tau(\phi_3) + \delta_1/2.
\end{eqnarray*}
With $\eps'$ and $\eps''$ sufficiently small, one has 
$$ \mathrm{d}_\tau((\tilde{d}_{1, \delta})^{\frac{1}{n}} d_{1, \delta}) > \mathrm{d}_\tau(\phi_3),\quad \tau \in \tr(A). $$
  

Note that $(\tilde{d}_{1, \delta})^{\frac{1}{n}} d_{1, \delta} \in \overline{\phi_2D\phi_2}$ (both $(\tilde{d}_{1, \delta})^{\frac{1}{n}}$ and  $d_{1, \delta}$ belong to $D$, so they commute). By Property (C), for any $\eps'''>0$ (to be determined later), there is $u_3 \in \overline{\phi_2 A \phi_2} + \Comp 1$ such that 
\begin{equation}\label{approx-C-3-h}
u_3\phi_3u_3^* \in _{\eps'''}^{\norm{\cdot}_2} \overline{(\tilde{d}_{1, \delta})^{\frac{1}{n}} d_{1, \delta} D (\tilde{d}_{1, \delta})^{\frac{1}{n}} d_{1, \delta}},
\end{equation}
$$\mathrm{dist}_{2, \tr(A)}(u_3 d u^*_3, D_1) < \eps''', \quad \mathrm{dist}_{2, \tr(A)}(u_3^*du_3, D_1) < \eps''',\quad d \in D_1, $$
and
\begin{equation}\label{approx-C-3-u}
\norm{u_3u_3^* - 1}_{2, \tr(A)},\quad  \norm{u_3^*u_3 - 1}_{2, \tr(A)} < \eps'''.
\end{equation}

Note that, by \eqref{pert-fur-hereditary},
$$\tilde{d}_{1, 1}c = c,\quad c \in \overline{(\tilde{d}_{1, \delta})^{\frac{1}{n}} d_{1, \delta} A (\tilde{d}_{1, \delta})^{\frac{1}{n}} d_{1, \delta} }.$$
By \eqref{approx-C-3-h}, there is $c \in \overline{(\tilde{d}_{1, \delta})^{\frac{1}{n}} d_{1, \delta}D (\tilde{d}_{1, \delta})^{\frac{1}{n}} d_{1, \delta}}$ such that $\norm{c} \leq 1$ and
$$ \norm{c - u_3\phi_3 u_3^*}_{2, \tr(A)} < \eps'''.$$ 
Then 
\begin{eqnarray}\label{comp-approx}
 ( u^*_2\eta_{\eps}(u_1^* \psi_1 u_1)u_2) (u_3 \phi_3 u_3^*) & \approx_{O(\eps')}^{\norm{\cdot}_2} &  ( u^*_2\eta_{\eps}(d_1)u_2) (u_3 \phi_3 u_3^*) \quad\quad ( \eqref{approx-C-1} )   \\
 & \approx_{O(\eps'')}^{\norm{\cdot}_2} & \eta_{\eps}(u_2^*d_1 u_2) (u_3 \phi_3 u_3^*) \quad\quad (\eqref{approx-C-2-u})   \nonumber \\
 & \approx_{O(\eps'')}^{\norm{\cdot}_2} & \eta_{\eps}(d_{1, 1}) (u_3 \phi_3 u_3^*) \quad\quad (\eqref{pert-C-step-2})   \nonumber \\
 & \approx_{\eps/2+O(\eps'')}^{\norm{\cdot}_2} &  \eta_{\eps}(\tilde{d}_{1, 1}) (u_3 \phi_3 u_3^*) \quad \quad (\eqref{pert-fur-hereditary} \eqref{pre-pert-delta}) \nonumber \\
 & \approx_{\eps'''}^{\norm{\cdot}_2} &  \eta_{\eps}(\tilde{d}_{1, 1})  c \nonumber \\
 &=&c \approx_{\eps'''}^{\norm{\cdot}_2} u_3 \phi_3 u_3^*. \nonumber
\end{eqnarray} 


Then, consider the contraction $$ u = u_1u_2u_3.$$
Since 
$$u_2 \in \overline{\phi_1 A \phi_1}+\Comp 1\quad \mathrm{and}\quad u_3 \in \overline{\phi_2 A \phi_2} + \Comp 1, $$ one has
$$(u_1u_2u_3)^*\psi_1(u_1u_2u_3) \in \overline{\phi_1 A \phi_1}.$$
Moreover,
\begin{eqnarray*}
(u_1u_2u_3)^*\eta_{\eps_0/2}(\psi_1)(u_1u_2u_3)  & \approx_{O(\eps')}^{\norm{\cdot}_2} & u^*_3u_2^*\eta_{\eps_0/2}(u_1^*\psi_1 u_1)u_2 u_3 \quad\quad (\eqref{approx-C-1-u}) \\
& \approx_{O(\eps')}^{\norm{\cdot}_2} & u^*_3(u_2^*\eta_{\eps_0/2}(d_1)u_2) u_3 \quad\quad (\eqref{approx-C-1}) \\
& \in_{\eps''}^{\norm{\cdot}_2} & \overline{\phi_2 A \phi_2}, \quad\quad (\eqref{approx-C-2-here})
\end{eqnarray*}
and 
\begin{eqnarray*}
\eta_{\eps}((u_1u_2u_3)^*\psi_1(u_1u_2u_3)) \phi_3 & \approx_{O(\eps' + \eps'')}^{\norm{\cdot}_2} & u_3^*u_2^*\eta_\eps(u_1^*\psi_1 u_1)u_2u_3 \phi_3 \quad\quad (\eqref{approx-C-1-u}, \eqref{approx-C-2-u}) \\
 & \approx_{\eps'''}^{\norm{\cdot}_2} & u^*_3 (u^*_2  \eta_{\eps}(u_1^* \psi_1 u_1)u_2)(u_3 \phi_3 u_3^*) u_3 \quad\quad (\eqref{approx-C-3-u}) \\
 & \approx_{\eps + O(\eps'+\eps''+\eps''')} & u_3^* (u_3\phi_3 u_3^*) u_3= \phi_3. \quad\quad(\eqref{comp-approx})
 \end{eqnarray*}
With $\eps', \eps'', \eps'''$ sufficiently small, the contraction $u$ has the desired property.
\end{proof}

\end{document}
